Small academic departments conducting machine learning research face a persistent infrastructure challenge. They cannot afford the dual-stack approach common in larger institutions, i.e., a traditional HPC scheduler such as Slurm for batch training jobs alongside a separate Kubernetes deployment for interactive services, yet operating without any orchestration leads to inefficient resource utilization, manual scheduling conflicts, and poor environment reproducibility. This thesis presents the design and implementation of a modular, Kubernetes-native compute cluster that consolidates batch workloads, interactive development environments, and supporting services onto a single lightweight control plane, purpose-built for resource-constrained academic settings.

The platform is built on K3s, a lightweight Kubernetes distribution whose minimal resource footprint makes a production-grade control plane feasible on a small cluster of consumer-grade machines. An intelligent scheduling layer based on Volcano provides queue-based fair sharing, gang scheduling for distributed jobs, and priority-driven preemption, compensating for the absence of hardware-level GPU partitioning (NVIDIA MIG) on consumer-grade GPUs. Complementary components include JupyterHub for interactive notebook access, Kube-Ray for distributed training, MLflow for experiment tracking, and a private container registry for reproducible environments.

The entire stack is designed for single-administrator operability: node bootstrapping is automated via Ansible, application lifecycle is managed through Helm charts deployed via CI/CD pipelines, and all configuration is version-controlled. The architecture is evaluated on a multi-node cluster of heterogeneous consumer-grade hardware (NVIDIA Titan and GeForce RTX GPUs, 12--24\,GB VRAM), demonstrating that a single, unified control plane can effectively serve the diverse workload requirements of an academic AI research lab without dedicated DevOps staff.
